% GENERATED FILE. Edit canonical lessons, then run npm run generate:books.
% canonical-source-hash: fnv1a64:fdda3517f6ab6420
% canonical-lessons: PT-C515-balneario, PT-C515-claque, PT-C515-lesao, PT-C515-lesionar, PT-C515-batota

\chapter{Lesões e batota}
\label{ch:pt-lesoes-e-batota}
\hlchaptermodality{515}

\begin{chapteropening}
\textbf{By the end of this chapter:} \emph{I can talk about injuries, cheating and the changing room.}

Complete the last lesson of chapter 515: I can talk about injuries, cheating and the changing room.
\end{chapteropening}

\section[o balneário]{o balneário — changing room}

\label{lesson:PT-C515-balneario}



\begin{quote}
Before the new one: say the Portuguese for judo, then the Portuguese for fencing.
\end{quote}

\subsection*{You'll want to know: o balneário}
\textbf{o balneário} — \textquotedblleft{}changing room\textquotedblright{}.

\begin{quote}
Depois do jogo, os jogadores tomam banho no balneário.
\end{quote}

\begin{quote}
\emph{After the match, the players shower in the changing room.}
\end{quote}

In Brazil: \emph{o vestiário}.

\begin{cousinweb}
\emph{balneário} is a learned word built on Latin \emph{balneārius} \textquotedblleft{}of a bath\textquotedblright{} (Lewis \& Short), from \emph{balneum} \textquotedblleft{}a bath\textquotedblright{}.
\end{cousinweb}

\begin{grammarlens}[title={What you've built}]
One more word for telling, explaining and describing at B1.
\end{grammarlens}

\subsection*{Your turn}
\begin{itemize}
\raggedright
  \item \emph{Say it:} \emph{o balneário}
  \item \emph{Say it:} \emph{o balneário}, once more
  \item \emph{Say it:} a sentence of your own with \emph{o balneário}
  \item \emph{Recall:} say the Portuguese for judo, then the Portuguese for fencing, then say \emph{o balneário} again
\end{itemize}

\begin{tcolorbox}[breakable,colback=teal!4,colframe=teal!35!black,title={Before you move on}]
What does o balneário mean? (\textquotedblleft{}Changing room\textquotedblright{}.) Say it once more.
\end{tcolorbox}
\section[a claque]{a claque — supporters' group, fans}

\label{lesson:PT-C515-claque}



\begin{quote}
Before the new one: say the Portuguese for fencing, then the Portuguese for changing room.
\end{quote}

\subsection*{You'll want to know: a claque}
\textbf{a claque} — \textquotedblleft{}supporters' group, fans\textquotedblright{}.

\begin{quote}
A claque cantou durante o jogo todo.
\end{quote}

\begin{quote}
\emph{The supporters sang for the whole match.}
\end{quote}

In Portugal, an organised group of fans.

\begin{cousinweb}
\emph{claque} is usually traced to French \emph{claque}, a group paid to applaud in a theatre, from \emph{claquer} \textquotedblleft{}to clap\textquotedblright{}, an imitative verb.
\end{cousinweb}

\begin{grammarlens}[title={What you've built}]
One more word for telling, explaining and describing at B1.
\end{grammarlens}

\subsection*{Your turn}
\begin{itemize}
\raggedright
  \item \emph{Say it:} \emph{a claque}
  \item \emph{Say it:} \emph{a claque}, once more
  \item \emph{Say it:} a sentence of your own with \emph{a claque}
  \item \emph{Recall:} say the Portuguese for fencing, then the Portuguese for changing room, then say \emph{a claque} again
\end{itemize}

\begin{tcolorbox}[breakable,colback=teal!4,colframe=teal!35!black,title={Before you move on}]
What does a claque mean? (\textquotedblleft{}Supporters' group, fans\textquotedblright{}.) Say it once more.
\end{tcolorbox}
\section[a lesão]{a lesão — injury}

\label{lesson:PT-C515-lesao}



\begin{quote}
Before the new one: say the Portuguese for changing room, then the Portuguese for supporters' group.
\end{quote}

\subsection*{You'll want to know: a lesão}
\textbf{a lesão} — \textquotedblleft{}injury\textquotedblright{}.

\begin{quote}
Uma lesão no joelho deixou-o fora da equipa.
\end{quote}

\begin{quote}
\emph{A knee injury kept him out of the team.}
\end{quote}

Plural \emph{lesões}.

\begin{cousinweb}
\emph{lesão} is a learned borrowing from Latin \emph{laesiō} \textquotedblleft{}an injury\textquotedblright{} (EtymDB), from \emph{laedere} \textquotedblleft{}to hurt\textquotedblright{}.
\end{cousinweb}

\begin{grammarlens}[title={What you've built}]
One more word for telling, explaining and describing at B1.
\end{grammarlens}

\subsection*{Your turn}
\begin{itemize}
\raggedright
  \item \emph{Say it:} \emph{a lesão}
  \item \emph{Say it:} \emph{a lesão}, once more
  \item \emph{Say it:} a sentence of your own with \emph{a lesão}
  \item \emph{Recall:} say the Portuguese for changing room, then the Portuguese for supporters' group, then say \emph{a lesão} again
\end{itemize}

\begin{tcolorbox}[breakable,colback=teal!4,colframe=teal!35!black,title={Before you move on}]
What does a lesão mean? (\textquotedblleft{}Injury\textquotedblright{}.) Say it once more.
\end{tcolorbox}
\section[lesionar]{lesionar — to get injured (lesionar-se); to injure}

\label{lesson:PT-C515-lesionar}



\begin{quote}
Before the new one: say the Portuguese for supporters' group, then the Portuguese for injury.
\end{quote}

\subsection*{You'll want to know: lesionar}
\textbf{lesionar} — \textquotedblleft{}to get injured (lesionar-se); to injure\textquotedblright{}.

\begin{quote}
O avançado lesionou-se no treino.
\end{quote}

\begin{quote}
\emph{The striker got injured in training.}
\end{quote}

Mostly reflexive, with the pronoun after the verb: \emph{lesionou-se} — he got injured. \emph{Estar lesionado} — to be injured.

\begin{cousinweb}
\emph{lesionar} is built on \emph{lesão} \textquotedblleft{}injury\textquotedblright{}, a learned borrowing from Latin \emph{laesiō} (EtymDB).
\end{cousinweb}

\begin{grammarlens}[title={What you've built}]
One more word for telling, explaining and describing at B1.
\end{grammarlens}

\subsection*{Your turn}
\begin{itemize}
\raggedright
  \item \emph{Say it:} \emph{lesionar}
  \item \emph{Say it:} \emph{lesionar}, once more
  \item \emph{Say it:} a sentence of your own with \emph{lesionar}
  \item \emph{Recall:} say the Portuguese for supporters' group, then the Portuguese for injury, then say \emph{lesionar} again
\end{itemize}

\begin{tcolorbox}[breakable,colback=teal!4,colframe=teal!35!black,title={Before you move on}]
What does lesionar mean? (\textquotedblleft{}To get injured (lesionar-se); to injure\textquotedblright{}.) Say it once more.
\end{tcolorbox}
\section[a batota]{a batota — cheating (at games)}

\label{lesson:PT-C515-batota}



\begin{quote}
Before the new one: say the Portuguese for injury, then the Portuguese for to get injured.
\end{quote}

\subsection*{You'll want to know: a batota}
\textbf{a batota} — \textquotedblleft{}cheating (at games)\textquotedblright{}.

\begin{quote}
O meu irmão faz batota nas cartas.
\end{quote}

\begin{quote}
\emph{My brother cheats at cards.}
\end{quote}

\emph{Fazer batota} — to cheat. European Portuguese; in Brazil, \emph{trapaça}.

\begin{cousinweb}
The origin of \emph{batota} is not settled.
\end{cousinweb}

\begin{grammarlens}[title={What you've built}]
One more word for telling, explaining and describing at B1.
\end{grammarlens}

\subsection*{Your turn}
\begin{itemize}
\raggedright
  \item \emph{Say it:} \emph{a batota}
  \item \emph{Say it:} \emph{a batota}, once more
  \item \emph{Say it:} a sentence of your own with \emph{a batota}
  \item \emph{Recall:} say the Portuguese for injury, then the Portuguese for to get injured, then say \emph{a batota} again
\end{itemize}

\begin{tcolorbox}[breakable,colback=teal!4,colframe=teal!35!black,title={Before you move on}]
What does a batota mean? (\textquotedblleft{}Cheating (at games)\textquotedblright{}.) Say it once more.
\end{tcolorbox}
